%==============================================================================
% CHAPTER 12: Operads for Designing Systems of Systems
%==============================================================================
\chapter{Operads for Designing Systems of Systems}
\label{ch:operads}
\cite{baez2021operads}

\begin{flushright}
\textit{John C. Baez, John D. Foley}\\
Communicated by \textit{Notices} Associate Editor Emilie Purvine
\end{flushright}

\epigraph{``Analogous to how a group abstracts the symmetries of an object, an operad abstracts operations that compose many objects into a single object.''}{--- Baez \& Foley} \cite{baez2021operads}

\begin{abstract}
``Systems of systems'' engineering composes independent systems (radar, comms, weapons, sensors) into a larger capability. This chapter introduces \textbf{operads}---a categorical tool that formalizes the \emph{syntax} of composition (\index{composition (operadic)}how things plug together) and their \emph{algebras} (the concrete semantics). Applied to network design, search-and-rescue planning, and task coordination, \cite{baez2021operads}.
\end{abstract}

%==============================================================================
\section{Introduction: The Composition Problem}

\noindent\textit{Mathematical notation follows standard conventions. Key terms are defined in the glossary.}\n
\label{sec:op:intro}

A smartphone is a system of systems: touchscreen, camera, CPU, radio, GPS, battery... The \emph{design} is a blueprint for how components connect. The \emph{behavior} depends on the specific hardware.

In defense (DARPA \index{DARPA CASCADE}CASCADE program), we need to compose:
\begin{itemize}
    \item Aircraft with range-limited comms
    \item Ships, helicopters, drones for search-and-rescue
    \item Sensors, effectors, decision-makers for task planning
\end{itemize}

\textbf{Challenge}: The composition rules (syntax) are different from the physical realization (semantics). Operads separate these, \cite{baez2021operads}.

%==============================================================================
\section{What Is an Operad?}
\label{sec:op:def}

\begin{definition}[Colored Operad]
An \textbf{operad} $\mathcal{O}$ consists of: \cite{baez2021operads}
\begin{itemize}
    \item A set of \textbf{types} (colors) $\mathrm{Ob}(\mathcal{O})$
    \item For each tuple $(X_1,\dots,X_n; Y)$ with $X_i, Y \in \mathrm{Ob}(\mathcal{O})$, a set of \textbf{operations} $\mathcal{O}(X_1,\dots,X_n; Y)$
    \item \textbf{Composition}: Given $f \in \mathcal{O}(X_1,\dots,X_n; Y)$ and $g_i \in \mathcal{O}(Z_{i1},\dots,Z_{ik_i}; X_i)$,
        \[
        f \circ (g_1,\dots,g_n) \in \mathcal{O}(Z_{11},\dots,Z_{nk_n}; Y)
        \]
    \item \textbf{Identities}: $\mathrm{id}_X \in \mathcal{O}(X; X)$
    \item \textbf{Symmetric group actions}: Permuting inputs
\end{itemize}
Satisfying associativity, identity, equivariance axioms.
\end{definition}

\begin{definition}[Operad Algebra]
An \textbf{algebra} of $\mathcal{O}$ assigns to each type $X$ a set $A(X)$ and to each operation $f \in \mathcal{O}(X_1,\dots,X_n; Y)$ a function
\[
A(f): A(X_1) \times \cdots \times A(X_n) \to A(Y)
\]
preserving composition, identities, and symmetries, \cite{baez2021operads}.
\end{definition}

\textbf{Intuition}: Operad = \emph{syntax} (how to compose); Algebra = \emph{semantics} (what the components actually are).

%==============================================================================
\section{Network Operads}
\label{sec:op:network}

\begin{example}[Aircraft Communication Network]
\textbf{Types}: Number of aircraft $n \in \mathbb{N}$.

\textbf{Operations}: $\mathcal{O}(n_1,\dots,n_k; m)$ = graphs with $n_1+\cdots+n_k$ input vertices and $m$ output vertices. An operation tells you \emph{which new edges to add} to connect the sub-networks.

\textbf{Algebra}: For each $n$, $A(n)$ = set of possible states of $n$ aircraft (positions, velocities, comms status). An operation acts by adding comms links between aircraft that are in range.
\end{example}

\begin{theorem}[Network Operad Construction]
For any class of networks (graphs, directed graphs, typed edges, etc.), there is a network operad whose operations are ``ways to connect networks by adding edges.'' \cite{baez2021operads}
\end{theorem}

%==============================================================================
\section{Application: Search and Rescue}
\label{sec:op:sar}

Problem: Distribute search assets (ships, planes, helicopters, small boats) to cover a large area after a disaster (Fastnet Race 1979, Sydney-Hobart 1998).

\begin{enumerate}
    \item \textbf{Design operad}: Assets have ``slots'' (helicopter carries 3 boats). Operations = ways to nest assets, \cite{baez2021operads}.
    \item \textbf{Algebra}: Actual assets at actual locations with fuel constraints.
    \item \textbf{Optimization}: Search over operad operations to maximize ``search effort'' (area covered per time) within budget, \cite{baez2021operads}.
\end{enumerate}

The operad structure lets you \emph{search the space of designs} algorithmically, \cite{baez2021operads}.

%==============================================================================
\section{From Design to Tasking}
\label{sec:op:tasking}

Same operad, different algebra: \cite{baez2021operads}
\begin{itemize}
    \item \textbf{Design algebra}: Assets at locations $\to$ network structure
    \item \textbf{Tasking algebra}: Primitive tasks (``fly to'', ``search'', ``rescue'') composed into mission plans
\end{itemize}

\textbf{Catalyst operads} (Baez et al.): Some operations enable other operations (a ``catalyst'' agent coordinates multiple agents). Used for multi-agent task plans, \cite{baez2021operads}.

Translation to \textbf{constraint programming} (MILP/CPLEX): Each operation $\to$ decision variables; composition $\to$ constraints, \cite{baez2021operads}.

%==============================================================================
\section{Levels of Abstraction}
\label{sec:op:abstraction}

A system design proceeds through levels:
\begin{enumerate}
    \item \textbf{Abstract}: ``Network of boats with comms''
    \item \textbf{Concrete}: ``3 RHIBs, 2 helicopters, SATCOM links''
    \item \textbf{Detailed}: Specific models, frequencies, protocols
\end{enumerate}

\textbf{Operad homomorphisms} and \textbf{algebra homomorphisms} formalize moving between levels. This is \emph{change of abstraction}.

%==============================================================================
\section{Code: Simple Operad in Python}
\label{sec:op:code}

\begin{lstlisting}[style=python,caption=Network operad sketch] \cite{baez2021operads}
from dataclasses import dataclass
from typing import List, Dict, Set, Callable
from abc import ABC, abstractmethod

@dataclass
class NetworkOperation:
    """Operation: connect input networks by adding edges"""
    input_types: List[int]  # sizes of input networks
    output_type: int        # size of output network
    edges: Set[tuple]       # (source_idx, target_idx) where idx is global
    
    def __call__(self, *networks):
        # Compose networks according to edges
        # networks[i] has size input_types[i]
        pass

class Operad:
    def __init__(self):
        self.operations: Dict[tuple, List[NetworkOperation]] = {}
    
    def add(self, op: NetworkOperation):
        key = (tuple(op.input_types), op.output_type)
        self.operations.setdefault(key, []).append(op)
    
    def compose(self, outer: NetworkOperation, inners: List[NetworkOperation]):
        # Build composed operation
        # Offset inner edge indices by input network sizes
        pass

class Algebra:
    def __init__(self, interpret_type: Callable[[int], set]):
        self.interpret_type = interpret_type  # n -> set of n-aircraft states
    
    def apply(self, op: NetworkOperation, *states):
        # Apply operation to concrete states
        pass

# Example: Aircraft comms
network_operad = Operad() \cite{baez2021operads}
# Add operations for connecting networks...

# Algebra: real aircraft positions and comms ranges
def aircraft_states(n):
    # Return set of possible states for n aircraft
    # Each state: positions, velocities, comms equipment
    pass

aircraft_algebra = Algebra(aircraft_states)
\end{lstlisting}

%==============================================================================
\section{Bridge to Chapter 13}
\label{sec:op:bridge}

Operads provide a \emph{static} compositional structure for system design. Chapter 13 applies \emph{dynamic} Bayesian inference to a different domain: astronomy. Both deal with \emph{hierarchical structure}---operads explicitly, astronomy through the hierarchy of data $\to$ model $\to$ cosmology, \cite{baez2021operads}.

%==============================================================================
\section{Further Reading}
\label{sec:op:refs}

\begin{itemize}
    \item Baez, Foley (2021). \textit{Operads for Designing Systems of Systems}. Notices, 68(6).
    \item Baez, Foley, Moeller, Pollard (2020). \textit{Network Models}. Theory Appl. Categ.
    \item Spivak (2013). \textit{The Operad of Wiring Diagrams}. arXiv:1305.0297.
    \item Spivak (2014). \textit{Category Theory for the Sciences}. MIT Press.
    \item Yau (2018). \textit{Operads of Wiring Diagrams}. Springer.
\end{itemize}

%==============================================================================
\section{Exercises}
%==============================================================================
% Solutions
%==============================================================================
\section{Solutions}
\label{sec:op:solutions}

\begin{solution}
\textbf{Exercise 1 (Operad composition):} The composition operation in an operad satisfies associativity: $(f \circ_i g) \circ_j h = f \circ_i (g \circ_j h)$. This ensures that composing operations in different orders gives the same result.
\end{solution}

\begin{solution}
\textbf{Exercise 2 (Network operad):} The network operad has operations corresponding to network topologies. Composition corresponds to substituting one network into a node of another.
\end{solution}

\begin{solution}
\textbf{Exercise 3 (DARPA CASCADE):} The CASCADE program uses operads to compose heterogeneous systems. Each system is an algebra over the operad, and composition is given by the operad action.
\end{solution}

\begin{solution}
\textbf{Exercise 4 (Research):} Operadic methods can be extended to stochastic systems and to systems with uncertainty. The compositional structure allows local verification to imply global properties.
\end{solution}

\label{sec:op:exercises}

\begin{exercise}
Define the operad of \emph{wiring diagrams} for digital circuits. Types = natural numbers (wire count). Operations = circuits with given input/output wires, \cite{baez2021operads}.
\end{exercise}

\begin{exercise}
For the aircraft comms operad, define an algebra where $A(n)$ is the set of possible communication graphs on $n$ aircraft with given range constraints, \cite{baez2021operads}.
\end{exercise}

\begin{exercise}
Design a ``catalyst'' operation for a search-and-rescue scenario where a helicopter coordinates the deployment of 3 small boats.
\end{exercise}

\begin{exercise}[Research]
Extend the network operad to include \emph{temporal} composition: operations that describe not just spatial connections but time-dependent interactions (e.g., intermittent connectivity), \cite{baez2021operads}.
\end{exercise}